\section{Introduction}\label{sec:rrm-introduction}

We use the chamber regulator with weights $(6,6\sqrt2,6\sqrt2,6)$
and Wilson parameter $0<\rho<1$. The operator $D^{\rm ch}[B]$
acts on four-component spinors with $48$ chamber labels per site,
antiperiodic in every coordinate; its argument $B$ is the bond phase.
The stencil is recalled in \cref{subsec:rrm-regulator}.

Fix the block construction of \cite[Proposition~4.1 and \S4.1--4.4]{QED3Integration},
with its positive charge ceilings $e_1\le e_0<1$ and the
common-pair theorem \cite[Theorem~D]{CoverTuning}.
Thus $L\ge13$ is odd, $M=L^m$ with $m\in\Z_{\ge2}$,
$0<\alpha<1$, $1-\alpha<\varepsilon\le1/200$, and the localization and cutoff
powers are $r_{\rm loc}=2$, $p_0=9$, $p=19$.
The fixed ceilings denoted by $e_0,e_1$ here are the
$\bar e_0,\bar e_1$ of \cite{QED3Integration};
they are distinct from the running charge $e_j$ at $j=0$.
These fixed constants include the photon and fermion inverse bounds,
the exact block step and the terminal conditions. The parameter input is Proposition~4.1 of \cite{QED3Integration},
which supplies the construction and charge ceilings before any
probe-dependent restriction. The response estimates are proved here.

For $0<\bar m\le1$ put
\[
 h=L^{\lfloor\log_L\bar m^{-1}\rfloor},\qquad
 \eta=L^{-N},\qquad q=e\sqrt\eta,\qquad K=N+\log_Lh.
\]
Let $r_K$ be the least nonnegative integer power of $L$ at least
$[-\log(e\sqrt h)]^2$, and set
\[
 R_{\rm top}=hMr_K,\qquad R^*=R_{\rm top}/L,\qquad
 n_*=R^*/\eta,\qquad \ell_c=R^*/(2L),\qquad N_*=8L^3.
\]
The finite threshold $N_{\rm ref}(e)$ and stopping level on $R^*$
are those of \cite[Proposition~6.1 and (53)]{CoverTuning}.
For $N\ge N_{\rm ref}(e)$ use the real pair
$P_N(e)=(m_{0,N},\varepsilon_{0,N}^{\rm v})$ selected on
$R_{\rm top}$ by \cite[Theorem~D]{CoverTuning}, and write
$\boldsymbol\mu_N=\bar m\eta+m_{0,N}$.
This same pair is used for every source and every twist.
Its vacuum factor cancels from the normalized response below.

Fix a real nonzero $\beta\in C_c^\infty(\R^3;\Lambda^2)$,
put $f=dd^*\beta$, and let $D_\beta$ be the side of a least
coordinate cube containing $\supp\beta$. Enlarge that cube by two
in each direction to obtain $Q_\beta$ of side
$D_\beta^{\rm buf}=D_\beta+4$ and center $z_\beta$.
With $J_\eta$ denoting cube averaging and $\delta_{\eta,\mu}$
the centered difference, define the covering bond cochain
\begin{equation}\label{eq:rrm-probe-definition}
 \widehat a^o_\nu(x)=
 -\eta^{1/2}\sum_{\mu\ne\nu}
 (\delta_{\eta,\mu}J_\eta\beta_{\mu\nu})
                         (\eta x+\eta e_\nu/2).
\end{equation}
A placement $v\in\eta\Z^3$ is admitted when
$v+Q_\beta\subset[\ell_c/8,7\ell_c/8]^3$.
Translate the probe by $v$, reflect it into the cells indexed by
$\mathcal C_*=(\Z/2L\Z)^3$, and periodize on $C(n_*)$.
The resulting cochains are $\widehat a_c$; the reflection maps
and their orientation action are specified in \cref{sec:reference-cells}.
They have mutually orthogonal curvatures, all of squared norm
$v_\eta=\|d\widehat a^o\|_{\ell^2}^2$.

Write $C(n_*)=\Z^3/n_*\Z^3$ with its oriented cubical cells.
Let $\gamma_{n_*}$ be the standard Gaussian on its exact real
two-cochains, with counting inner product. For such a
curvature $u$, let $A(u)$ be its mean-zero transverse potential.
Set $H_{\theta,\nu}=\theta_\nu/(qn_*)$ and
$a(t)=e^{-1}\sum_ct_c\widehat a_c$. Define
\begin{equation}\label{eq:rrm-normalized-response}
 \mathring\Phi(t)=
 \frac{\displaystyle\int_{[-\pi,\pi]^3}\frac{\dd\theta}{(2\pi)^3}
       \int\det\bigl(D^{\rm ch}[q(A(u)+H_\theta+a(t))]
                         +\boldsymbol\mu_N\bigr)\,\dd\gamma_{n_*}(u)}
      {\displaystyle\int_{[-\pi,\pi]^3}\frac{\dd\theta}{(2\pi)^3}
       \int\det\bigl(D^{\rm ch}[q(A(u)+H_\theta)]
                         +\boldsymbol\mu_N\bigr)\,\dd\gamma_{n_*}(u)}.
\end{equation}
The denominator is positive, including when the tuned total mass
has either sign. The free one-cell ratio uses the positive reference mass:
\begin{equation}\label{eq:response-free-ratio}
 N_0(t)=\frac{\det\bigl(D^{\rm ch}[q(t/e)\widehat a_0]
                                      +\bar m\eta\bigr)}
                   {\det\bigl(D^{\rm ch}[0]+\bar m\eta\bigr)},
 \qquad c_m=N_0^{(m)}(0),\quad m=2,4.
\end{equation}
These are ordinary derivatives of the ratio. Put
\[
 \Theta_m^{(0)}=
 \left|\left[\prod_{c\in\mathcal C_*}
           (\partial_{t_c}^m-c_m)\right]\mathring\Phi(0)
                                      \right|^{1/N_*},\qquad
 \mathfrak e=\tfrac12(1/4-9\varepsilon)-12\varepsilon>0.
\]

\begin{theorem}[Normalized reference response]\label{thm:rrm-main}
For the fixed construction, mass $0<\bar m\le1$, and probe $\beta$
above, there is $e_{\rm RR}>0$ such that for every
$0<e\le e_{\rm RR}$ and every integer $N\ge N_{\rm ref}(e)$,
the function \eqref{eq:rrm-normalized-response} is entire and
\begin{equation}\label{eq:rrm-main-bound}
 \Theta_4^{(0)}\le C e^{\mathfrak e},\qquad
 (1+|c_2|)\Theta_2^{(0)}\le C e^{\mathfrak e}.
\end{equation}
The constants are independent of $N$ and of the admitted placement.
Every centered placement,
$|v+z_\beta-(\ell_c/2)(1,1,1)|_\infty\le\eta/2$, is admitted
in this range. One may take
$C=3\cdot10^5(1+\sqrt h\,C_\beta)^4h^{1/8}$,
where the finite, charge-independent $C_\beta$ is defined in
\eqref{eq:interacting-uniform-probe-constants}.

Let $\mathbb E^{(0)}$ denote the normalized zero-flux curvature law
on this same reference torus and at the same tuned mass.
Equivalently it is the probability law of $F^{\rm u}=u$ obtained
by weighting $\gamma_{n_*}$ with the zero-source twist-averaged
determinant in \eqref{eq:rrm-normalized-response}.
For $X_c=\langle d\widehat a_c,F^{\rm u}\rangle$ and
$H_2(y)=y^2-v_\eta$, $H_4(y)=y^4-6v_\eta y^2+3v_\eta^2$,
\begin{equation}\label{eq:rrm-main-wick}
 \mathbb E^{(0)}\prod_{c\in\mathcal C_*}
             (e^{-m}H_m(X_c)-c_m)
       =\left[\prod_{c\in\mathcal C_*}
           (\partial_{t_c}^m-c_m)\right]\mathring\Phi(0),
 \qquad m=2,4.
\end{equation}
All expectations are finite. The response estimate concerns the full
fixed family of $N_*=8L^3$ cells and the harmonic-averaged zero-flux law.
\end{theorem}

The proof has two distinct parts. The exact identities require no
renormalization estimates. Integer cochain unfolding retains the gauge
and harmonic Jacobians; a real source shift, followed by analytic
continuation, then gives \eqref{eq:rrm-main-wick}.
For the quantitative bound, the terminal density retains the external
connection throughout the common trajectory. Its logarithm is separated
into the free one-cell logarithms, their errors, and terms involving
several cells. A tree expansion preserves the distance between marked
cells and bounds the latter terms exponentially in their separation.

The restricted terminal density defines a positive measure of mass at
most one. The omitted configurations have small probability under the
zero-flux Gaussian bound. A two-constants estimate transfers this
smallness to the complex coefficient norm needed for differentiation.
In the ring retaining degree at most four in each source, every error
family which misses a cell is annihilated by that cell's free factor.
The surviving families therefore carry a power of the error for every
cell. This proves the estimate after taking the $N_*$th root.
The final section verifies all additional response inequalities on one
positive charge interval, including at jumps of the rounded scales.

These are the source and response arguments of
\cite[\S19--20, \S21.1, \S23]{QED3Master}.
The terminal operator estimates and exact Borel density used below are
proved in \cite[Theorem~B, \S6 and Appendix~A]{TerminalPressure}.
The free ratio has the convention of
\cite[Theorem~B]{PositiveProbe}; its coefficient estimates at the
analytic radius required here are proved in \cref{subsec:interacting-one-cell}.
No sign condition on the continuum fourth coefficient is needed.
Transfer to larger tori and other flux sectors is a subsequent result.

\subsection{Regulator and elementary conventions}\label{subsec:rrm-regulator}
Let $\sigma_\mu$ be the Pauli matrices and
$\gamma_\mu=\sigma_3\otimes\sigma_\mu$.
The labels $g\in O_{\rm cub}=O(3)\cap\mathrm{GL}(3,\Z)$ are
the signed permutation matrices. Let $s_0$ reverse coordinate three,
$s_1,s_2$ exchange $(2,3)$ and $(1,2)$, and let $r_{e_1}$ reverse
coordinate one. Put
\[
 n_0=e_3,\quad n_1=(e_2-e_3)/\sqrt2,\quad
 n_2=(e_1-e_2)/\sqrt2,\quad n_3=-e_1,\qquad
 M_i(g)=\tfrac12w_i(\rho+\gamma\cdot gn_i).
\]
On a unit torus the chamber operator is
\begin{equation}\label{eq:rrm-regulator}
 \begin{split}
 (D^{\rm ch}[B]\psi)(c,g)={}&
 \sum_{i=0}^2M_i(g)\bigl(\psi(c,g)-\psi(c,gs_i)\bigr)\\
 &+M_3(g)\bigl(\psi(c,g)
       -\mathrm e^{iB(c,c+ge_1)}\psi(c+ge_1,gr_{e_1})\bigr).
 \end{split}
\end{equation}
Reversed bonds carry the negative cochain value.
We write $D_A=D^{\rm ch}[qA]$ in the unit-lattice source identities.
At mesh $\eta_j$, the external hop and the bond phase have the
corresponding factors $\eta_j^{-1}$ and $e_j\eta_j$.
The spinor pairing has weight $1/48$ per chamber label; its constant
weight does not change a finite-dimensional determinant.

For smooth functions,
\[
 J_\eta v(y)=\eta^{-3}\int_{[-\eta/2,\eta/2]^3}v(y+z)\,\dd z,
 \qquad
 \delta_{\eta,\mu}v(y)=
 \frac{v(y+\eta e_\mu/2)-v(y-\eta e_\mu/2)}{\eta}.
\]
We use $\tau_vf(x)=f(x-v)$ and the seminorm
$|\beta|_{C^j}=\max_{\mu<\nu,|\alpha|\le j}
\|\partial^\alpha\beta_{\mu\nu}\|_\infty$.
Counting cochain pairings are bilinear after complexification.
Distances are periodic sup distances. The symbol $e$ is the physical
charge; numerical exponentials are written $\exp(\cdot)$ or $\mathrm e$.
